% kotexgweb.tex -- Korean localization and a LuaTeX PDF back end for GWEB.
%
% Put  \input kotexgweb.tex  in the limbo of a .w file and process the woven
% output with luatex (not pdftex):  luatex foo.tex.  gweave emits
% "\input gwebmac" as the very first line, so this file is read afterwards and
% overrides gwebmac wherever Korean text or the LuaTeX engine needs it.
%
% Requires luatexko and the Noto Serif/Sans CJK KR fonts. The Hangul typefaces
% are set in hangulfont.tex (\input below); edit that file alone to change them.

% ---- Korean fonts -----------------------------------------------------------
% The Hangul typefaces -- and their luatexko settings -- live in hangulfont.tex,
% so switching to a different Korean font is a change to that one file and never
% touches this one. It defines the body fonts \tenmj ... \twelvemj, the
% math-hangul fonts, the title companions \hangultitlefont/\hantitlefont, and
% \sectionfont, all used below.
\input hangulfont

% size-switch macros (borrowed from manmac.tex) that bind each CM family to its
% matching hangul font, so \rm/\it/\bf/\tt carry Korean as well as Latin.
\catcode`@=11 % borrow the private macros of PLAIN (with care)

\font\tentex=cmtex10

\font\inchhigh=cminch
% Contents-page title (\con) and the typewriter limbo title: pair gwebmac's Latin
% title fonts with the Hangul companions \hangultitlefont/\hantitlefont from
% hangulfont.tex, so only Korean characters take the CJK face -- the same
% "\latinfont\hangulfont" idiom used for \rm above.
\let\latintitlefont=\titlefont
\def\titlefont{\latintitlefont\hangultitlefont}

\let\lattitlefont=\ttitlefont
\def\ttitlefont{\lattitlefont\hantitlefont}

\font\ninerm=cmr9
\let\mc=\ninerm % medium caps
\font\eightrm=cmr8
\font\sixrm=cmr6

\font\ninei=cmmi9
\font\eighti=cmmi8
\font\sixi=cmmi6
\skewchar\ninei='177 \skewchar\eighti='177 \skewchar\sixi='177

\font\ninesy=cmsy9
\font\eightsy=cmsy8
\font\sixsy=cmsy6
\skewchar\ninesy='60 \skewchar\eightsy='60 \skewchar\sixsy='60

\font\eightss=cmssq8

\font\eightssi=cmssqi8

\font\ninebf=cmbx9
\font\eightbf=cmbx8
\font\sixbf=cmbx6

\font\ninett=cmtt9
\font\eighttt=cmtt8

\hyphenchar\tentt=-1 % inhibit hyphenation in typewriter type
\hyphenchar\ninett=-1
\hyphenchar\eighttt=-1

\font\ninesl=cmsl9
\font\eightsl=cmsl8

\font\nineit=cmti9
\font\eightit=cmti8

\font\tenu=cmu10 % unslanted text italic
\font\magnifiedfiverm=cmr5 at 10pt
\font\cmman=cmman % font used for miscellaneous Computer Modern variations

%
\newskip\ttglue
\def\tenpoint{\def\rm{\fam0\tenrm\tenmj}%
  \textfont0=\tenrm \scriptfont0=\sevenrm \scriptscriptfont0=\fiverm
  \textfont1=\teni \scriptfont1=\seveni \scriptscriptfont1=\fivei
  \textfont2=\tensy \scriptfont2=\sevensy \scriptscriptfont2=\fivesy
  \textfont3=\tenex \scriptfont3=\tenex \scriptscriptfont3=\tenex
  \setmathhangulfonts\texthangul\scripthangul\scriptscripthangul
  \def\it{\fam\itfam\tenit\tengtd}%
  \textfont\itfam=\tenit
  \def\sl{\fam\slfam\tensl\tengt}%
  \textfont\slfam=\tensl
  \def\bf{\fam\bffam\tenbf\tenbd}%
  \textfont\bffam=\tenbf \scriptfont\bffam=\sevenbf
  \scriptscriptfont\bffam=\fivebf
  \def\tt{\fam\ttfam\tentt\tentz}%
  \textfont\ttfam=\tentt
  \tt \ttglue=.5em plus.25em minus.15em
  \normalbaselineskip=14pt
  \let\sc=\eightrm
  \let\big=\tenbig
  \setbox\strutbox=\hbox{\vrule height8.5pt depth3.5pt width\z@}%
  \normalbaselines\rm}

\def\ninepoint{\def\rm{\fam0\ninerm\ninemj}%
  \textfont0=\ninerm \scriptfont0=\sixrm \scriptscriptfont0=\fiverm
  \textfont1=\ninei \scriptfont1=\sixi \scriptscriptfont1=\fivei
  \textfont2=\ninesy \scriptfont2=\sixsy \scriptscriptfont2=\fivesy
  \textfont3=\tenex \scriptfont3=\tenex \scriptscriptfont3=\tenex
  \setmathhangulfonts\texthangulnine\scripthangulsix\scriptscripthangul
  \def\it{\fam\itfam\nineit\ninegtd}%
  \textfont\itfam=\nineit
  \def\sl{\fam\slfam\ninesl\ninegt}%
  \textfont\slfam=\ninesl
  \def\bf{\fam\bffam\ninebf\ninebd}%
  \textfont\bffam=\ninebf \scriptfont\bffam=\sixbf
   \scriptscriptfont\bffam=\fivebf
  \def\tt{\fam\ttfam\ninett\ninetz}%
  \textfont\ttfam=\ninett
  \tt \ttglue=.5em plus.25em minus.15em
  \normalbaselineskip=13pt
  \let\sc=\sevenrm
  \let\big=\ninebig
  \setbox\strutbox=\hbox{\vrule height8pt depth3pt width\z@}%
  \normalbaselines\rm}

\def\eightpoint{\def\rm{\fam0\eightrm\eightmj}%
  \textfont0=\eightrm \scriptfont0=\sixrm \scriptscriptfont0=\fiverm
  \textfont1=\eighti \scriptfont1=\sixi \scriptscriptfont1=\fivei
  \textfont2=\eightsy \scriptfont2=\sixsy \scriptscriptfont2=\fivesy
  \textfont3=\tenex \scriptfont3=\tenex \scriptscriptfont3=\tenex
  \setmathhangulfonts\texthanguleight\scripthangulsix\scriptscripthangul
  \def\it{\fam\itfam\eightit\eightgtd}%
  \textfont\itfam=\eightit
  \def\sl{\fam\slfam\eightsl\eightgt}%
  \textfont\slfam=\eightsl
  \def\bf{\fam\bffam\eightbf\eightbd}%
  \textfont\bffam=\eightbf \scriptfont\bffam=\sixbf
   \scriptscriptfont\bffam=\fivebf
  \def\tt{\fam\ttfam\eighttt\eighttz}%
  \textfont\ttfam=\eighttt
  \tt \ttglue=.5em plus.25em minus.15em
  \normalbaselineskip=11pt
  \let\sc=\sixrm
  \let\big=\eightbig
  \setbox\strutbox=\hbox{\vrule height7pt depth2pt width\z@}%
  \normalbaselines\rm}

\def\tenmath{\tenpoint\fam-1 } % use after $ in ninepoint sections
\def\tenbig#1{{\hbox{$\left#1\vbox to8.5pt{}\right.\n@space$}}}
\def\ninebig#1{{\hbox{$\textfont0=\tenrm\textfont2=\tensy
  \left#1\vbox to7.25pt{}\right.\n@space$}}}
\def\eightbig#1{{\hbox{$\textfont0=\ninerm\textfont2=\ninesy
  \left#1\vbox to6.5pt{}\right.\n@space$}}}

\catcode`\@=12
\tenpoint

\fontdimen7\tentex=0pt  % restore cmtex10's "no extra space after sentence"
                        % (the \font\tentex above drops gwebmac's setting)

% Displayed code sets at \codebaselineskip (gwebmac's default is 12pt, plain TeX's
% leading for the Latin 10pt code font). Korean prose is set at 15pt
% (\normalbaselineskip above), so a bit more air suits the code beside it; 13.5pt
% loosens it a touch while staying well under the prose's 15pt.
\codebaselineskip=13.5pt

% gwebmac sets strings (\ST) and @d names (\MAC) in cmtex10, which has no
% Hangul; pair it with the monospace hangul (\tentz) so Korean inside a Go
% string or an @d name is also typewriter, not the serif fallback.
\def\ST#1{\hbox{\tentex\tentz #1}}
\def\MAC#1{\hbox{\tentex\tentz #1}}

% \.{...} (typewriter text in comments and the limbo) is cmtex10 too; bring the
% monospace hangul (\tentz) into scope so Korean inside \.{...} is typewriter as
% well. The original \. then sets cmtex10 for the Latin part as before.
% \leavevmode has to stay outside the group.  A paragraph that begins with
% \.{...} would otherwise begin inside it, so \everypar would fire inside the
% group; a hook that clears itself -- \everypar{\hskip-\parindent\everypar{}},
% the usual way of unindenting the paragraph after a heading -- would then be
% restored when the group closed and would go on unindenting later paragraphs.
% (gwebmac's own \. opens a group before any horizontal material too, so the
% \leavevmode has to come ahead of both groups.)
\let\ttdot\.
\def\.{\leavevmode\kottdot}
\def\kottdot#1{{\tentz\ttdot{#1}}}

% The small print (cross-reference notes, contents columns, section-name notes)
% is cmr8, which has no Hangul; pair it with the 8pt serif hangul (\eightmj) so
% Korean there is small too, instead of falling back to the body size. (The
% running head's \headfont stays sans, set earlier.)
\def\smallfont{\eightrm\eightmj}

% \datethis: localize the timestamp. cweb sets this line in \eightrm (8pt); pair
% the 8pt serif hangul (\eightmj) so the Korean date is the same 8pt as cweb's,
% not the serif body size. \today reads "2026년 6월 29일"; \dateat is a space.
\def\today{{\number\year}년 {\number\month}월 {\number\day}일}
\def\datefont{\eightrm\eightmj}
\def\dateat{\ }

% ---- Korean wording (the GWEB analogue of cwebman's localization list) ------
% Cross-reference notes printed beneath a code definition. Korean has no
% singular/plural inflection, so the "section"/"sections" pair reads the same.
\def\U#1{\note{이 코드는 #1절에서 쓰인다}}
\def\Us#1{\note{이 코드는 #1절에서 쓰인다}}
\def\A#1{\note{이 코드는 #1절에도 정의되어 있다}}
\def\As#1{\note{이 코드는 #1절에도 정의되어 있다}}
% The change-file roster printed just before the index.
\def\ch#1{\note{변경 파일이 바꾼 절: #1}}
% The usage note beside an entry in the section-name list.
\def\Nused#1{#1절에서 쓰임}
\def\Nuseds#1{#1절에서 쓰임}
% The running head over the list of section names, and the contents columns.
\def\namesheading{절 이름 목록}
\def\outsecname{절 이름 목록} % title of the section-names PDF bookmark
\def\sectionword{절}
\def\pageword{쪽}
% Running-head font: pair the small roman (cmr8) with a sans hangul (\ninegt,
% Noto Sans) so Korean in the header matches the upper-cased Latin instead of
% clashing with it as the serif body font would. This is the group (chapter)
% title's font, and \groupfont follows it.
\def\headfont{\eightrm\ninegt}
% The document title, though, is set as written (not upper-cased), and on the
% contents page it is serif; so when the author supplies a \def\title we set it
% in serif hangul (\ninemj, Noto Serif) here too, to match. A default job-name
% title keeps the sans of the group titles, since it reads like one of them.
\def\Gtitlefont{\headfont}

% ---- LuaTeX PDF back end: blue links and Korean (UTF-16BE) bookmarks --------
% gwebmac's \dest/\link/\bookmark are written for pdfTeX primitives that
% luatex does not provide; we restate them with luatex's \pdfextension forms.
% A PDF outline title must be a UTF-16BE string, so Korean titles are converted
% by a small Lua helper, gweb_pdfutf16(s), which prints s as the body of a PDF
% literal string: the UTF-16BE byte-order mark followed by each code unit as a
% pair of octal-escaped bytes (surrogate-paired above the BMP).
%
% The helper is built inline. Two cautions about Lua inside \directlua: (1) the
% TeX-special characters % # _ & ^ ~ $ are made catcode 12 below (% is Lua's
% modulo operator, not a TeX comment; # is its length operator), and the
% backslash is produced as string.char(92) so none appears in the source; and
% (2) \directlua turns the source's newlines into spaces, so a Lua "--" comment
% would run to the end of the whole chunk -- hence there are none here.
\begingroup
\catcode`\%=12 \catcode`\#=12 \catcode`\_=12 \catcode`\&=12
\catcode`\^=12 \catcode`\~=12 \catcode`\$=12
\directlua{
  local function oct(v)
    local d1 = math.floor(v / 64); v = v % 64
    local d2 = math.floor(v / 8); local d3 = v % 8
    return tostring(d1) .. tostring(d2) .. tostring(d3)
  end
  function gweb_pdfutf16(s)
    local bs = string.char(92)
    local t = { bs .. oct(254) .. bs .. oct(255) }
    for _, c in utf8.codes(s) do
      if c >= 65536 then
        c = c - 65536
        local w1 = 55296 + math.floor(c / 1024)
        local w2 = 56320 + (c % 1024)
        t[#t + 1] = bs .. oct(math.floor(w1 / 256)) .. bs .. oct(w1 % 256)
        t[#t + 1] = bs .. oct(math.floor(w2 / 256)) .. bs .. oct(w2 % 256)
      else
        t[#t + 1] = bs .. oct(math.floor(c / 256)) .. bs .. oct(c % 256)
      end
    end
    tex.sprint(-2, table.concat(t))
  end
}
\endgroup
\def\dest#1{\pdfextension dest num #1 fith\relax}
\def\link#1#2{\pdfextension startlink attr{/Border[0 0 0]}goto num #1\relax
  \pdfextension literal{0 0 1 rg}#2\pdfextension literal{0 g}\pdfextension endlink}
\def\bookmark#1#2#3#4{\pdfextension outline goto num #2 count #3
  {\directlua{gweb_pdfutf16("\luaescapestring{#4}")}}}
% gwebmac's \pdfURL is likewise written for pdfTeX primitives; restate it (the
% link text may of course be Korean).
\def\pdfURL#1#2{%
  \pdfextension startlink attr{/Border[0 0 0]}
    user{/Subtype /Link /A << /S /URI /URI (#2) >>}\relax
  \pdfextension literal{0 0 1 rg}#1\pdfextension literal{0 g}%
  \pdfextension endlink}
% Open the bookmark (outline) pane when the document is first shown.
\pdfextension catalog{/PageMode /UseOutlines}
